% !TeX program = xelatex
% Final exam cheat sheet: 2 A4 sheets, printed double-sided (4 pages), no colour.
% Topics follow the lecturer's exam scope. Build: xelatex final-cheat-sheet.tex
\documentclass[10pt]{extarticle}
\usepackage[a4paper,margin=0.28in,includefoot,footskip=0.18in]{geometry}
\usepackage{fontspec}
\usepackage{microtype}
\usepackage{amsmath,amssymb}
\usepackage{multicol}
\usepackage{enumitem}
\usepackage{booktabs}
\usepackage{xcolor}
% array pinned: the local texmf array.sty needs a newer kernel (see README)
\usepackage{array}[=2024-06-01]
\usepackage{tabularx}
\usepackage{needspace}

\pagestyle{plain}
\setlength{\columnsep}{0.14in}
\setlength{\columnseprule}{0.3pt}
\setlength{\parindent}{0pt}
\setlength{\parskip}{1pt}
\setlist{nosep,leftmargin=0.95em,labelsep=0.3em}
\setlength{\tabcolsep}{2.5pt}
\renewcommand{\arraystretch}{1.02}
\AtBeginDocument{%
  \setlength{\abovedisplayskip}{1.5pt}\setlength{\belowdisplayskip}{1.5pt}%
  \setlength{\abovedisplayshortskip}{1pt}\setlength{\belowdisplayshortskip}{1pt}}

% topic banner and numbered block heading
\newcommand{\topic}[1]{\par\needspace{8\baselineskip}\vspace{3pt}\noindent
  \colorbox{black!12}{\parbox{\dimexpr\linewidth-2\fboxsep}{\bfseries\small #1}}\par\vspace{1pt}}
\newcounter{cs}
\newcommand{\cs}[1]{\par\needspace{5\baselineskip}\vspace{2.5pt}\refstepcounter{cs}%
  \noindent{\bfseries\thecs.\ #1}\par\vspace{0.5pt}\hrule\vspace{1.5pt}}
\newcommand{\key}[1]{\par\smallskip\noindent\fbox{\parbox{\dimexpr\linewidth-2\fboxsep-2\fboxrule}{\centering$\displaystyle #1$}}\par\smallskip}
\newcommand{\CPI}{\mathrm{CPI}}
\newcommand{\IC}{\mathrm{IC}}
\newcommand{\ET}{\mathrm{ET}}
\newcommand{\MR}{\mathrm{MR}}
\newcommand{\MP}{\mathrm{MP}}
\newcommand{\HT}{\mathrm{HT}}
\newcommand{\AMAT}{\mathrm{AMAT}}
\newcommand{\tw}[1]{\texttt{#1}}

\begin{document}
\begin{center}
  {\large\bfseries CO5059 Advanced System Architectures -- Final Cheat Sheet}\\[0.5pt]
  {\scriptsize Personal study aid; H\&P \emph{Computer Architecture: A Quantitative
  Approach} 5th ed.\ and the course slides.}
\end{center}
\vspace{-3pt}

\raggedcolumns
\begin{multicols*}{3}
\raggedright

% #############################################################################
\topic{1 \quad PERFORMANCE}

\cs{Performance and speedup}
$\text{Perf}=1/\ET$; $\text{Speedup}_{A/B}=\ET_B/\ET_A=\text{old}/\text{new}$.
\begin{itemize}
  \item ``X is $n\%$ faster than Y'': $\ET_Y/\ET_X=1+n/100$; divide by the
  \emph{faster} time (15 s vs.\ 10 s $=50\%$ faster).
  \item CPU time $=$ user $+$ system; \%CPU $=(u+s)$/elapsed (can exceed 100\%
  with several cores).
  \item Latency $=$ time per task; throughput $=$ tasks/time; one task at a time:
  throughput $=1/$latency.
\end{itemize}

\cs{CPU equation and CPI}
\key{T=\IC\times\CPI\times C=\frac{\IC\times\CPI}{f}}
$\text{Cycles}=\sum\CPI_i\IC_i$;\quad $\CPI=\sum\CPI_iF_i$,\ $F_i=\IC_i/\IC$.\\
Time fraction of class $i=\CPI_iF_i/\CPI$.
\begin{itemize}
  \item Program/compiler/ISA $\to$ IC, CPI; organization $\to$ CPI, $C$;
  technology $\to C$.
  \item Compare total cycles, not IC or CPI alone (removing cheap instructions
  raises CPI but saves time).
  \item $\text{MIPS}=\IC/(T\cdot10^6)=f/(\CPI\cdot10^6)$; fair only for the same
  ISA, program, compiler. $\text{MFLOPS}=\#\text{FP}/(T\cdot10^6)$.
  \item 200 MHz $\to$ 5 ns; 2.5 GHz $\to$ 0.4 ns; 3 GHz $\to$ 0.333 ns.
\end{itemize}
\textbf{Ex.} ALU 50\%/1, Load 20\%/5, Store 10\%/3, Br 20\%/2: $\CPI=2.2$; load time
share $1.0/2.2=45\%$.
\textbf{Ex.} 4M instr, 2.5 GHz, 8 ms $\Rightarrow\CPI=8\,\text{ms}\cdot2.5\,\text{GHz}/4\text{M}=5$;
new compiler 5M instr, CPI 2: 4 ms, speedup 2.

\cs{CPI with stalls}
$\CPI=\CPI_{\text{ideal}}+$ stalls/instr; each source $=$ frequency $\times$ fraction
penalized $\times$ penalty. Memory: $+\frac{\text{acc}}{\text{instr}}\times\MR\times\MP$
(acc $=1+f_{ld/st}$).

\cs{Amdahl's law}
$F$ = fraction of \emph{old time} enhanced, $S$ = its speedup.
\key{\text{Speedup}=\frac{1}{(1-F)+F/S}\le\frac{1}{1-F}}
\begin{itemize}
  \item $\ET_{new}=\ET_{old}[(1-F)+F/S]$.
  \item Several disjoint: $1/\big((1-\sum F_i)+\sum F_i/S_i\big)$.
  \item Parallel: $1/((1-p)+p/N)$.
  \item Target $T$: $S=\dfrac{F}{1/T-(1-F)}$; impossible if $T\ge\frac1{1-F}$.
  \item From a mix: $F=\CPI_iF_i/\CPI$, \emph{not} $F_i$.
  \item $f'$ measured on the \emph{new} time: Speedup $=(1-f')+f'S$,
  $F=f'S/((1-f')+f'S)$.
  \item Clock slowed by $k$: Speedup $=1/\big([(1-F)+F/S]\,k\big)$.
\end{itemize}
\textbf{Ex.} mul 60\% ($\times6$), div 30\% ($\times12$): mul alone $1/(0.4+0.1)=2$;
div alone $1.38$ ($\le1/0.7$); both $1/0.225=4.44$.

% #############################################################################
\topic{2 \quad INSTRUCTION SET ARCHITECTURE}

\cs{ISA styles ($C=A+B$)}
{\scriptsize
\begin{tabularx}{\linewidth}{@{}lX@{}}
\toprule
Stack & \tw{PUSH A; PUSH B; ADD; POP C}; dense code, stack bottleneck \\
Accum. & \tw{LOAD A; ADD B; STORE C}; one register AC, heavy memory traffic \\
Mem-mem & \tw{ADD C,A,B}; all operands in memory \\
Reg-mem & \tw{LOAD R1,A; ADD R3,R1,B; STORE R3,C} (x86, 68000, VAX) \\
Load-store & \tw{LOAD R1,A; LOAD R2,B; ADD R3,R1,R2; STORE R3,C}; MIPS; easy to pipeline \\
\bottomrule
\end{tabularx}}
Design: operations, addressing modes, data types, format, registers.
Instruction types: data processing, data storage (LD/ST), data movement (MOV),
control. Cycle: IF, ID, operand fetch, EX, WB.

\cs{Addressing modes}
{\scriptsize
\begin{tabularx}{\linewidth}{@{}lX@{}}
Immediate & \tw{ADD R0,\#10} \\
Register & \tw{ADD R0,R1} \\
Direct & \tw{ADD R0,(100)}: M[100] \\
Reg.\ indirect & \tw{ADD R0,(R1)}: M[R1] \\
Displacement & \tw{LD R1,100(R2)}: M[R2+100] \\
Indexed & \tw{(R1+R2)}: M[R1+R2] \\
Mem.\ indirect & M[M[R2]] \\
PC-relative & PC $+$ offset \\
Scaled & M[100+R3+R4$\cdot d$] \\
\end{tabularx}}
Displacement and immediate dominate (SPEC89); most displacements fit in 12--16
bits; small constants $\approx50\%$ of operands.

\cs{Format, CISC vs.\ RISC}
\begin{itemize}
  \item Variable length (VAX, x86): dense, hard decode; fixed (MIPS, ARM): easy
  decode and pipelining; hybrid (MIPS16, Thumb).
  \item Field for $n$ choices: $\lceil\log_2n\rceil$ bits; $k$-bit signed:
  $-2^{k-1}\ldots2^{k-1}-1$. Expanding opcode: spare codes become escapes.
  \item More registers: fewer spills, but more bits per instruction, slower
  register file, more state to save.
  \item CISC: fewer, complex, variable-length instructions; many modes (68020:
  25); costly memory. RISC: fixed 32-bit, load-store, few modes, CPI$\to$1, more
  instructions.
\end{itemize}

\cs{MIPS}
{\scriptsize
\begin{tabular*}{\linewidth}{@{\extracolsep{\fill}}lcccccc@{}}
\toprule
& 31--26 & 25--21 & 20--16 & 15--11 & 10--6 & 5--0\\
\midrule
R & op 0 & rs & rt & rd & shamt & funct\\
I & op & rs & rt & \multicolumn{3}{c}{imm 16}\\
J & op & \multicolumn{5}{c}{address 26}\\
\bottomrule
\end{tabular*}}
\begin{itemize}
  \item Branch: target $=$ PC$+4+$SignExt(imm)$\times4$; offset
  $=(T-(\text{PC}+4))/4$; reach $\pm2^{15}$ instructions.
  \item J: $\{(\text{PC}+4)[31{:}28],\text{addr},00\}$, same 256 MiB; field $=T/4$.
  \tw{jr}: any address.
  \item Shifts: source \tw{rt}, \tw{rs}$=0$. I-format destination \tw{rt}.
  \item Sign-extend: \tw{addi slti}, ld/st offsets, branches. Zero-extend:
  \tw{andi ori xori}.
  \item \tw{add addi sub} trap on overflow; \tw{addu addiu subu} never. \tw{slt}
  signed, \tw{sltu} unsigned. \tw{lb lh} sign-extend, \tw{lbu lhu} zero-extend;
  words aligned.
  \item No MOV (\tw{add \$d,\$s,\$zero}), NOT (\tw{nor \$d,\$s,\$zero}), SUBI
  (\tw{addi} $-k$), SLA ($=$ \tw{sll}).
  \item 32-bit constant: \tw{lui t,hi; ori t,t,lo} (\tw{addi} breaks if lo $\ge$
  \tw{0x8000}); \tw{lui 0x1234; addi 0x5678} $=$ \tw{0x12345678}.
  \item \tw{mult/div} $\to$ HI, LO (LO quotient, HI remainder).
  \item Big-endian: MSB at lowest address (course); x86 little-endian (MSB at
  highest).
  \item \tw{ll/sc}: \tw{sc} returns 1 if no write since \tw{ll}, else 0: retry.
  \item \tw{jal}: \tw{\$ra}$\leftarrow$PC$+4$; non-leaf saves \tw{\$ra}; stack
  grows down. Preserved: \tw{\$s0--7 \$sp \$fp \$ra \$gp}; not: \tw{\$t \$a \$v}.
\end{itemize}
{\scriptsize
\textbf{funct} (op 0): sll 00, srl 02, sra 03, jr 08, jalr 09, add 20, addu 21,
sub 22, subu 23, and 24, or 25, xor 26, nor 27, slt 2a, sltu 2b.\\
\textbf{opcode}: j 02, jal 03, beq 04, bne 05, addi 08, addiu 09, slti 0a, sltiu 0b,
andi 0c, ori 0d, xori 0e, lui 0f, lb 20, lh 21, lw 23, lbu 24, lhu 25, sb 28, sh 29,
sw 2b, ll 30, sc 38 (hex).\par}
\$zero 0, \$at 1, \$v0--1 2--3, \$a0--3 4--7, \$t0--7 8--15, \$s0--7 16--23,
\$t8--9 24--25, \$gp 28, \$sp 29, \$fp 30, \$ra 31.

% #############################################################################
\topic{3 \quad PIPELINING}

\cs{Timing}
\key{\text{Cycles}=k+(n-1)+\text{stalls},\quad t_{clk}=\max t_i+t_{latch}}
\begin{itemize}
  \item Unpipelined: $\sum t_i$ per instruction. Ideal speedup
  $\frac{nk}{k+n-1}\to k$.
  \item $\CPI=1+$stalls; Speedup $=\frac{\CPI_u}{\CPI_p}\cdot\frac{C_u}{C_p}
  =\frac{\text{depth}}{1+\text{stalls}}$.
  \item 100 tasks, 6 stages, 2 ns: $(6+99)\cdot2=210$ ns.
  \item No common clock (lecturer): $T(n)=\sum t_i+(n-1)\max t_i$. Stages
  5,4,5,10,4: 3 instructions $=28+20=48$ ns; instruction $i$ waits $(i-1)(10-5)$:
  the 4th waits 15 ns.
  \item Improves throughput, not latency; unbalanced stages hurt.
\end{itemize}

\cs{Stages and hazards}
IF, ID (read registers), EX, MEM, WB. Register file: write first half, read second
half. Hazards: structural (one memory port: split I/D), data, control. Dependences:
RAW (true), WAR (anti), WAW (output); \textbf{RAR is not a hazard}. 5-stage: only
RAW stalls.
{\scriptsize
\begin{tabularx}{\linewidth}{@{}Xc@{}}
\toprule
Producer $\to$ consumer at distance $d$ & stalls\\
\midrule
no forwarding & $\max(0,3-d)$\\
fwd ALU$\to$ALU (EX/MEM, MEM/WB$\to$EX) & 0\\
fwd load$\to$use, $d=1$ & 1\\
fwd load$\to$use $d\ge2$; load$\to$store data & 0\\
\bottomrule
\end{tabularx}}
Diagram: a stalled instruction waits in ID (\textit{st} after ID); the next waits
in IF; fetch when IF is free. Total $=$ last WB. Reordering must keep RAW, WAR and
WAW (rename if needed). Forwarding is hardware; scheduling, delay slots and static
prediction are compiler techniques.

\cs{Lecturer's pipeline problem}
\tw{lw t0; lw t1; add t3,t0,t1; sw t3; lw t0; add t3,t0,t1; sw t3}: no fwd
\textbf{19} cycles (8 stalls); fwd \textbf{13} (2 load-use stalls); reorder
(1)(2)(3)(5)(4)(6)(7): \textbf{12} (moving (5) above (3) is a WAR bug); with
renaming 11.

\cs{Control hazards}
Penalty $=$ resolving stage $-1$ (MEM: 3, ID: 1).
$\CPI=1+f_{br}\times\text{fraction penalized}\times\text{penalty}$.
Stall/flush; predict not taken (taken branches pay); predict taken (no gain in
the 5-stage MIPS); delayed branch: slot from before (best), from the target (useful
if taken), from the fall-through (useful if not). $\approx60\%$ of slots filled,
80\% of those usefully.
\textbf{Ex.} 20\% branches, 60\% taken: MEM stall 1.6; MEM predict-not-taken 1.36;
ID stall 1.2; ID predict-not-taken 1.12.

% #############################################################################
\topic{4 \quad CACHE}

\cs{Address breakdown}
$C$ size, $B$ block, $a$ ways, $A$ address bits.
\key{S=\frac{C}{B\cdot a},\ O=\log_2B,\ I=\log_2S,\ \text{tag}=A-I-O}
Block addr $=\lfloor\text{addr}/B\rfloor$; set $=$ block mod $S$; tag
$=\lfloor\text{block}/S\rfloor$. DM: $a=1$; FA: $S=1$, no index. Doubling $a$:
index $-1$, tag $+1$. Bits per block $=8B+$tag$+$V($+$D)($+$LRU).
\textbf{Ex.} 8 KB 2-way 16 B, 32-bit: O 4, I 8, tag 20; addresses 0x00 0x14 0x08 0x10
0x20 $\to$ blocks 0 1 0 1 2: M M H H M $=$ \textbf{3 misses}.
\textbf{Ex.} Byte 1200, DM 64 blocks of 16 B: block 75, line 11.

\cs{Four questions}
Placement (DM / set associative / FA); identification (tag $+$ valid);
replacement; write strategy.

\cs{Replacement}
\begin{itemize}
  \item Random; LRU (update on every access; $\lceil\log_2 n!\rceil$ bits); FIFO
  (oldest; hits do not reorder).
  \item \textbf{Tree-LRU}: $n-1$ bits; 4-way: $b_0$ chooses the pair, $b_1$ in
  \{0,1\}, $b_2$ in \{2,3\}; an access sets the path bits to point \emph{away};
  victim: follow the bits. w0: $b_0{=}1,b_1{=}1$; w1: $b_0{=}1,b_1{=}0$; w2:
  $b_0{=}0,b_2{=}1$; w3: $b_0{=}0,b_2{=}0$.
  \item \textbf{Clock}: reference bit per way $+$ pointer; an access sets the bit;
  to replace: bit 1 $\to$ clear and advance; bit 0 $\to$ evict. New block gets bit
  1; the pointer moves past it.
  \item H\&P pseudo-LRU: one bit per way, all but the newest cleared when full.
  \item Replacement has a second-order effect (only on misses).
\end{itemize}

\cs{Write policies}
{\scriptsize
\begin{tabularx}{\linewidth}{@{}lX@{}}
\toprule
Write-through & cache $+$ lower level; no dirty bit; write buffer; memory current \\
Write-back & cache only; dirty bit; written on eviction; less traffic \\
Write allocate & write miss: fetch the block, then write \\
No-write alloc. & write miss: write the lower level only \\
\bottomrule
\end{tabularx}}
Pairs: WB$+$WA, WT$+$NWA. W100 W100 R200 W200 W100: NWA 4 misses 1 hit; WA 2 misses 3 hits.

\cs{The 3 Cs}
Compulsory (first access; infinite cache); capacity (also misses in an FA cache of
the same size); conflict (hits in FA, misses here); $+$ coherence. 8-way $\approx$
FA; 2:1 rule: DM of size $N$ $\approx$ 2-way of size $N/2$.

\cs{Performance}
\key{\AMAT=\HT+\MR\times\MP}
\[\text{CPU time}=\IC\Big(\CPI_{ex}+\tfrac{\text{acc}}{\text{instr}}\MR\cdot\MP\Big)C\]
Misses/instr $=\MR\times$ acc/instr. Split I/D: \%I $=1/(1+f_{data})$; a unified
cache adds a cycle per data access (one port).
\textbf{Ex.} $\CPI_{ex}=1$, 1.5 acc, 2\%, MP 25: $\CPI=1.75$.
\textbf{Ex.} Hit 1, miss 105, no cache 100: 5\% $\to$ 6.2 (or 6.25); random accesses
$\to\approx105$: without locality a cache hurts.

\cs{Multilevel}
\key{\AMAT=\HT_1+\MR_1(\HT_2+\MR_{2}\cdot\MP_2)}
Local $\MR_2=$ misses$_2$/accesses$_2$ (accesses$_2=$ misses$_1$); global $=\MR_1\MR_2$.
Stalls/instr $=\frac{\text{m}_1}{\text{I}}\HT_2+\frac{\text{m}_2}{\text{I}}\MP_2$.
L2 reduces the L1 \emph{miss penalty}. Inclusion: L1 $\subset$ L2; exclusion: swap.
\textbf{Ex.} 1000 refs, 40 L1, 20 L2 misses; 1/10/200 cycles: local 50\%, global
2\%, $\AMAT=1+0.04(10+0.5\cdot200)=5.4$. \textbf{Ex.} 1 ns, 10\%, 100 ns $\to$ 11 ns;
add L2 8 ns, 20\%, 105 ns $\to$ 3.9 ns.

\cs{Optimizations}
{\scriptsize
\begin{tabularx}{\linewidth}{@{}lX@{}}
\toprule
Larger block & $\downarrow$compulsory; $\uparrow$MP, conflict \\
Larger cache & $\downarrow$capacity; $\uparrow$HT, cost \\
Higher assoc. & $\downarrow$conflict; $\uparrow$HT \\
Multilevel & $\downarrow$MP \\
Reads before writes & $\downarrow$MP; check the write buffer (RAW) \\
Avoid translation & $\downarrow$HT: VIPT, $C/a\le$ page size \\
\bottomrule
\end{tabularx}}
Advanced: small simple L1, way prediction, pipelined access, nonblocking,
multibanked (bank $=$ block mod $n$), critical word first, merging write buffer,
compiler (blocking), HW/SW prefetching.
\textbf{Ex.} L1 A: 2-way 0.6 ns 5\%; B: 8-way 0.8 ns 4\%: B wins only if MP $>20$ ns.

% #############################################################################
\topic{5 \quad VIRTUAL MEMORY}

\cs{Basics}
Main memory caches the disk; private space per process; program $>$ RAM, isolation,
relocation, sharing, non-contiguous pages. Page $=$ frame size (4 KiB). PTE: PFN,
valid/present, R/W/X, U/S, dirty, reference. The page table (indexed by VPN) stores
\textbf{physical page numbers}. Placement fully associative; replacement by the OS
(LRU approx., use bit); writes always \textbf{write-back}. Internal fragmentation
(paging) vs.\ external (segmentation).

\cs{Translation}
\key{\text{VPN}=\lfloor\text{VA}/P\rfloor,\ \text{off}=\text{VA}\bmod P,\ \text{PA}=\text{PFN}\cdot P+\text{off}}
4 KiB: low 3 hex digits $=$ offset: \tw{0x00403A7C}, VPN 403$\to$1F2: \tw{0x1F2A7C}.
VA 3500, $P=1024$, page 3$\to$7: offset 428, PA 7596.
\begin{itemize}
  \item Flat table $=2^{\text{VPN bits}}\times$ PTE, per process: 32-bit, 4 KiB, 4 B
  $\to$ 4 MiB (100 processes: 400 MiB).
  \item Multilevel: field $=\log_2(P/\text{PTE})$ bits: $10+10+12$; x86-64
  $4\times9+12$; a walk costs $k$ accesses.
  \item Inverted: one entry per frame, hashed on the VPN.
\end{itemize}

\cs{TLB}
Small (fully associative) cache of VPN$\to$PFN. Hit: no table access. Miss: walk;
page present $\to$ refill, else \textbf{page fault}. TLB miss $\ne$ page fault (may
or may not occur).
\key{\text{EAT}=h(t+m)+(1-h)(t+km+m)}
$t=20$, $m=100$, $h=98\%$: 122 ns. Reach $=$ entries $\times$ page. Context switch:
flush or ASID. Impossible: TLB hit with page fault; page fault with cache hit.

\cs{Page faults}
Trap $\to$ check legality $\to$ free frame or victim (write back if dirty) $\to$
load $\to$ update PTE $\to$ restart. $\text{EAT}=(1-p)m+p\,s$ (slide:
$T_{mem}+(1-h)T_{disk}$). 100 ns, 10 ms, 0.1\%: $\approx10{,}100$ ns. Thrashing:
working sets $>$ RAM.

\cs{Page replacement}
OPT (next use furthest; lower bound); FIFO (oldest; Belady: \tw{123412512345} 9
faults/3 frames, 10/4); LRU (a stack algorithm); LFU; Clock.
\tw{121132143112415621}, 3 frames: FIFO \textbf{9}, LRU \textbf{11}, OPT 7, Clock 9.
\tw{70120304}: FIFO 7, LRU 6. With $k$ empty frames the first $k$ distinct pages fault.

\cs{Segmentation and protection}
(seg, off): if off $<$ limit then PA $=$ base $+$ off, else trap. (2,53), base 4300,
limit 400 $\to$ 4353; (2,400) traps. Paging vs.\ segments: fixed/variable size, 1/2
address words, internal/external fragmentation, trivial/hard replacement. Paged
segments combine both. x86 real mode: seg$\times16+$off. Protection: isolation,
R/W/X, user/kernel; needs $\ge2$ modes, state users cannot write, system call and
return.

\cs{VA to cache (VIPT)}
Index within the page offset: $I+O\le$ offset bits $\iff C/a\le P$.
64-bit VA, 41-bit PA, 8 KB pages: VPN 51 $=$ TLB tag 43 $+$ index 8; L1 index 7 $+$
offset 6; L1 tag 28; L2 $19+16+6$.
{\scriptsize
\begin{tabular*}{\linewidth}{@{\extracolsep{\fill}}lll@{}}
\toprule
& L1 cache & VM\\
\midrule
Unit & 16--128 B & 4--64 KiB\\
Hit & 1--3 cycles & 100--200 cycles\\
Miss penalty & 8--200 cycles & $10^6$--$10^7$ cycles\\
Miss rate & 0.1--10\% & $10^{-5}$--$10^{-3}$\%\\
Replacement & hardware & OS\\
\bottomrule
\end{tabular*}}

% #############################################################################
\topic{6 \quad MULTIPROCESSORS}

\cs{Classification}
Flynn: SISD (uniprocessor), SIMD (vector, GPU), \textbf{MISD (none built)}, MIMD
(multicore). UMA/SMP: one shared memory, equal latency, bus limits scaling.
NUMA/DSM: distributed memory, one address space, remote accesses slower. Message
passing: clusters. Walls: power, ILP, memory. Speedup 80 on 100 processors needs
99.75\% parallel. Remote 666 cycles at 0.2\%, base CPI 0.5: $\CPI=1.7$.

\cs{Synchronization}
Race: shared data $+$ a write $+$ no synchronization. Critical section: mutual
exclusion, progress, bounded waiting. Atomic primitives: test-and-set, exchange,
fetch-and-increment, CAS, \tw{ll/sc}.
\textbf{Spin lock}: loop on exchange (each try is a write $\to$ invalidations).
\textbf{TTAS}: spin on a cached read; exchange only when free. Backoff; queue (MCS)
locks. Mutex: sleeps; for long waits. Deadlock: circular wait $\to$ one lock order.
\textbf{Barrier}: counter $+$ release flag; reuse can deadlock $\to$
sense-reversing; tree barrier at scale.

\cs{Consistency}
Coherence: one location. Consistency: order across locations. \textbf{SC}: one global
order, each thread in program order. \tw{x=1; r0=y} $\|$ \tw{y=1; r1=x}: $r0=r1=0$
impossible under SC, possible with a store buffer. TSO/PC: W$\to$R relaxed (x86);
PSO: also W$\to$W; weak ordering: only syncs ordered; release consistency:
acquire/release.

\cs{Dependences and ILP}
True (RAW, real); anti (WAR) and output (WAW) are name dependences, removed by
\textbf{renaming}; control; resource. Bernstein: $P_1\|P_2$ iff
$I_1\cap O_2=I_2\cap O_1=O_1\cap O_2=\emptyset$. GCD test: store \tw{X[ai+b]}, load
\tw{X[ci+d]}: a dependence needs $\gcd(a,c)\mid(d-b)$. Loop-carried: DOALL /
DOACROSS. Unrolling \tw{x[i]+=s}: 9 $\to$ 7 cycles; unrolled $\times4$: 3.5/element.
\begin{itemize}
  \item \textbf{Superpipeline}: more, shorter stages, faster clock, 1 issue per
  cycle, larger penalties. 12 stages, 2 ns, 100 instr: 18.5 ns (5 stages 41.6, none
  200).
  \item \textbf{Superscalar}: $n$ issues per cycle, hardware checks (IPC $>1$).
  \item In-order issue and completion: RAW. In-order issue, out-of-order
  completion: $+$WAW (ROB). Out-of-order issue: $+$WAR (renaming, window, ROB).
  Tomasulo renames; speculation commits in order.
  \item \textbf{VLIW}: the compiler packs operations (unrolling, software
  pipelining, trace scheduling); simple hardware, low power; nops increase code
  size; no binary compatibility.
\end{itemize}

\cs{Multithreading}
Per-thread PC and registers. \textbf{Coarse}: switch on long stalls (refill
cost). \textbf{Fine}: switch every cycle; hides short stalls, slows each thread.
\textbf{SMT}: issue from several threads \emph{in the same cycle}: TLP $+$ ILP;
fills horizontal waste.

\cs{Coherence}
Coherent: (1) a processor reads back its own write, (2) others' writes are seen
eventually, (3) writes are serialized. Gives migration and replication.
\textbf{Snooping}: broadcast bus, every cache watches; small SMPs.
\textbf{Directory}: per block a state (U/S/M) $+$ sharer bit vector at the home
node; point to point; scales. Write invalidate (used) vs.\ write update.
{\scriptsize
\begin{tabularx}{\linewidth}{@{}lX@{}}
\toprule
Read miss & I$\to$S (MESI: E if no other copy) \\
Write & I/S$\to$M; invalidate other copies \\
Write in E & E$\to$M \textbf{silent} (no bus) \\
Snooped read & M$\to$S (flush), E$\to$S \\
Snooped write & M (flush), E, S $\to$ I \\
Evict M & write back \\
\bottomrule
\end{tabularx}}
\textbf{Purpose of E}: avoid a second bus transaction when a thread writes a line
it read. Read-then-write of private $Y$: MSI 2 transactions, MESI 1. False sharing:
different words in one block $\to$ pad.
\textbf{Review answers}: renaming removes WAR and WAW; SMT $=$ TLP $+$ ILP; MISD is
unused; VLIW relies on the compiler.

\cs{Exam checklist}
\begin{itemize}
  \item ns vs.\ cycles; ms $\to\times10^6$ ns; bytes vs.\ words vs.\ bits.
  \item $F$ in Amdahl is a fraction of \emph{time}.
  \item Count the instruction fetch in accesses/instr; local vs.\ global $\MR$.
  \item Branch offset counts words from PC$+4$.
  \item State the initial cache/TLB/frame contents in traces.
\end{itemize}

\end{multicols*}

\clearpage
\noindent\colorbox{black!12}{\parbox{\dimexpr\linewidth-2\fboxsep}{\bfseries\small WORKED TRACES (lecturer's review and course examples)}}

\medskip
\noindent\textbf{A. Pipeline diagrams} (5-stage, write first half / read second half; \textit{st} $=$ STALL; stalled instruction waits in ID, the one behind waits in IF).
Code: \tw{(1) lw \$t0,0(\$t2); (2) lw \$t1,4(\$t4); (3) add \$t3,\$t0,\$t1; (4) sw \$t3,8(\$t2); (5) lw \$t0,12(\$t4); (6) add \$t3,\$t0,\$t1; (7) sw \$t3,16(\$t2)}.

\noindent\textbf{No forwarding: 19 cycles, 8 stalls} (consumer's ID in the cycle of the producer's WB or later).
\begin{center}\footnotesize\setlength{\tabcolsep}{2.4pt}
\begin{tabular}{@{}lccccccccccccccccccc@{}}
\toprule
Instruction & 1 & 2 & 3 & 4 & 5 & 6 & 7 & 8 & 9 & 10 & 11 & 12 & 13 & 14 & 15 & 16 & 17 & 18 & 19 \\
\midrule
(1) \texttt{lw \$t0,0(\$t2)} & IF & ID & EX & MEM & WB &  &  &  &  &  &  &  &  &  &  &  &  &  &  \\
(2) \texttt{lw \$t1,4(\$t4)} &  & IF & ID & EX & MEM & WB &  &  &  &  &  &  &  &  &  &  &  &  &  \\
(3) \texttt{add \$t3,\$t0,\$t1} &  &  & IF & ID & \textit{st} & \textit{st} & EX & MEM & WB &  &  &  &  &  &  &  &  &  &  \\
(4) \texttt{sw \$t3,8(\$t2)} &  &  &  & IF & \textit{st} & \textit{st} & ID & \textit{st} & \textit{st} & EX & MEM & WB &  &  &  &  &  &  &  \\
(5) \texttt{lw \$t0,12(\$t4)} &  &  &  &  &  &  & IF & \textit{st} & \textit{st} & ID & EX & MEM & WB &  &  &  &  &  &  \\
(6) \texttt{add \$t3,\$t0,\$t1} &  &  &  &  &  &  &  &  &  & IF & ID & \textit{st} & \textit{st} & EX & MEM & WB &  &  &  \\
(7) \texttt{sw \$t3,16(\$t2)} &  &  &  &  &  &  &  &  &  &  & IF & \textit{st} & \textit{st} & ID & \textit{st} & \textit{st} & EX & MEM & WB \\
\bottomrule
\end{tabular}
\end{center}
\noindent\textbf{Forwarding: 13 cycles, 2 stalls} (only load-use at distance 1). Best legal reordering (1)(2)(3)(5)(4)(6)(7): 12 cycles; with renaming of (5)--(7): 11.
\begin{center}\footnotesize\setlength{\tabcolsep}{2.4pt}
\begin{tabular}{@{}lccccccccccccc@{}}
\toprule
Instruction & 1 & 2 & 3 & 4 & 5 & 6 & 7 & 8 & 9 & 10 & 11 & 12 & 13 \\
\midrule
(1) \texttt{lw \$t0,0(\$t2)} & IF & ID & EX & MEM & WB &  &  &  &  &  &  &  &  \\
(2) \texttt{lw \$t1,4(\$t4)} &  & IF & ID & EX & MEM & WB &  &  &  &  &  &  &  \\
(3) \texttt{add \$t3,\$t0,\$t1} &  &  & IF & ID & \textit{st} & EX & MEM & WB &  &  &  &  &  \\
(4) \texttt{sw \$t3,8(\$t2)} &  &  &  & IF & \textit{st} & ID & EX & MEM & WB &  &  &  &  \\
(5) \texttt{lw \$t0,12(\$t4)} &  &  &  &  &  & IF & ID & EX & MEM & WB &  &  &  \\
(6) \texttt{add \$t3,\$t0,\$t1} &  &  &  &  &  &  & IF & ID & \textit{st} & EX & MEM & WB &  \\
(7) \texttt{sw \$t3,16(\$t2)} &  &  &  &  &  &  &  & IF & \textit{st} & ID & EX & MEM & WB \\
\bottomrule
\end{tabular}
\end{center}
\begin{multicols}{2}
\noindent\textbf{B. One 4-way set}, refs \tw{0 1 2 3 0 4 1 0 5 2 3 0} (fill ways 0--3; tree bits $b_0b_1b_2$ from 000; Clock gives new blocks bit 1). M(victim), H.

{\footnotesize\setlength{\tabcolsep}{1.6pt}
\begin{tabular}{@{}lcccccccccccc@{}}
\toprule
Ref & 0 & 1 & 2 & 3 & 0 & 4 & 1 & 0 & 5 & 2 & 3 & 0 \\
\midrule
LRU & M & M & M & M & H & M1 & M2 & H & M3 & M4 & M1 & H \\
FIFO & M & M & M & M & H & M0 & H & M1 & M2 & M3 & M4 & H \\
Tree & M & M & M & M & H & M2 & H & H & M3 & M1 & M4 & H \\
bits & 110 & 100 & 001 & 000 & 110 & 011 & 101 & 111 & 010 & 100 & 001 & 111 \\
Clock & M & M & M & M & H & M0 & H & M2 & M3 & M1 & M4 & H \\
\bottomrule
\end{tabular}}

LRU 9, FIFO 9, Tree 8, Clock 9 misses.

\medskip
\noindent\textbf{C. Page replacement}, 3 frames, \tw{1 2 1 1 3 2 1 4 3 1 1 2 4 1 5 6 2 1}; F(victim).

{\footnotesize\setlength{\tabcolsep}{1.3pt}
\begin{tabular}{@{}l*{18}{c}@{}}
\toprule
& 1&2&1&1&3&2&1&4&3&1&1&2&4&1&5&6&2&1\\
\midrule
FIFO&F&F&&&F&&&F1&&F2&&F3&&&F4&F1&&F2\\
LRU&F&F&&&F&&&F3&F2&&&F4&F3&&F2&F4&F1&F5\\
OPT&F&F&&&F&&&F2&&&&F3&&&F4&F5&&\\
\bottomrule
\end{tabular}}

FIFO 9, LRU 11, OPT 7 (Clock 9).

\columnbreak
\noindent\textbf{D. Coherence trace}, 3 CPUs, $X$ uncached; states P1 P2 P3.

{\footnotesize
\begin{tabular}{@{}lllll@{}}
\toprule
Access & MSI & MSI bus & MESI & MESI bus \\
\midrule
P1 rd & S I I & BusRd & E I I & BusRd \\
P2 rd & S S I & BusRd & S S I & BusRd \\
P1 wr & M I I & invalidate & M I I & invalidate \\
P3 rd & S I S & BusRd, P1 flush & S I S & same \\
P2 wr & I M I & write miss & I M I & same \\
P1 rd & S S I & BusRd, P2 flush & S S I & same \\
\bottomrule
\end{tabular}}

Private $Y$ read then written: MSI BusRd $+$ invalidate (2); MESI BusRd (E) then silent E$\to$M (1).

\medskip
\noindent\textbf{E. MIPS encodings} (fields; hex)\\
{\footnotesize
\tw{add \$t0,\$s1,\$s2}: 0 17 18 8 0 0x20 $=$ \tw{0x02324020}\\
\tw{slt \$t0,\$s1,\$s2}: 0 17 18 8 0 0x2a $=$ \tw{0x0232402A}\\
\tw{srl \$t1,\$s3,4}: 0 0 19 9 4 0x02 $=$ \tw{0x00134902}\\
\tw{lw \$t0,16(\$t1)}: 0x23 9 8 16 $=$ \tw{0x8D280010}\\
\tw{sw \$ra,4(\$sp)}: 0x2b 29 31 4 $=$ \tw{0xAFBF0004}\\
\tw{beq \$s0,\$s1,256}: 4 16 17 256 $=$ \tw{0x12110100}\\
\tw{jal 0x00400200}: 3, addr \tw{0x0100080} $=$ \tw{0x0C100080}\\
\tw{0x1520FFFC} $=$ \tw{bne \$t1,\$zero,-4}: target PC$+4-16$.}
\end{multicols}
\end{document}
